% ═══════════════════════════════════════════════════════════════
%  chapters/ch10-dac-dma.tex
%  DMA circular и калибровка генератора
% ═══════════════════════════════════════════════════════════════

\chapter{DMA circular и калибровка}
\label{ch:dac-dma}

\epigraph{<<DMA~--- это как конвейер. Он работает, пока ты
занимаешься другими делами>>.}{}

\section{Введение}
\label{sec:dac-dma-intro}

В \texttt{ch09} мы написали \emph{таблицы} форм и
\emph{DDS} \emph{аккумулятор}. Но \emph{есть} \emph{проблема}:
если \emph{вызывать} \texttt{dds\_tick()} \emph{из}
\emph{прерывания}, \emph{каждый} \emph{такт} \emph{CPU}
\emph{тратит} \emph{время}. \emph{Для} \SI{1}{\mega\hertz}~---
\emph{миллион} \emph{прерываний} \emph{в} \emph{секунду}.
\emph{CPU} \emph{не} \emph{успевает} \emph{ничего} \emph{другого}.

\textbf{Решение}: \textbf{DMA circular}. \emph{DMA} \emph{сам}
\emph{читает} \emph{таблицу} \emph{и} \emph{пишет} \emph{в} DAC.
\emph{CPU} \emph{свободен}.

\begin{ideabox}
\textbf{Аналогия: автопилот.} \emph{Без} DMA~--- \emph{пилот}
(CPU) \emph{сам} \emph{крутит} \emph{штурвал} \emph{каждую}
\emph{секунду}. \emph{С} DMA~--- \emph{автопилот} \emph{крутит}
\emph{штурвал}, \emph{а} \emph{пилот} \emph{отдыхает} \emph{или}
\emph{решает} \emph{другие} \emph{задачи}.
\end{ideabox}

\subsection{Что вы узнаете}

\begin{itemize}
    \item Как работает \textbf{DMA circular}.
    \item Почему \emph{CPU} \emph{свободен}.
    \item Как \emph{откалибровать} сигнал.
    \item Как \emph{усилить} сигнал.
    \item Как \emph{реализовать} это на STM32.
\end{itemize}

\section{DMA circular}
\label{sec:dac-dma-what}

\textbf{DMA} \emph{circular}~--- это \emph{режим}, при
котором DMA \emph{читает} \emph{буфер} \emph{по} \emph{кругу}.
Когда \emph{дошёл} \emph{до} \emph{конца}~--- \emph{начинает}
\emph{с} \emph{начала}.

\begin{figure}[htbp]
\centering
\adjustbox{max width=\textwidth}{%
\begin{tikzpicture}[
    x=1cm, y=1cm,
    cell/.style={
        draw, thick,
        minimum width=1.2cm,
        minimum height=1.2cm,
        align=center,
        font=\footnotesize\ttfamily
    },
    full/.style={cell, fill=clrLogicData!30},
    active/.style={cell, fill=clrWarn!30},
    arr/.style={-{Latex[length=3mm]}, very thick}
]

% ─── Координатная сетка (для отладки) ────────────
\draw[step=1cm, gray!20, very thin] (-1, -1) grid (8, 3);

% Ячейки
\node[full]   (c0) at (0, 0) {A};
\node[full]   (c1) at (1.3, 0) {B};
\node[active] (c2) at (2.6, 0) {C};
\node[full]   (c3) at (3.9, 0) {D};
\node[full]   (c4) at (5.2, 0) {E};
\node[full]   (c5) at (6.5, 0) {F};

% Стрелки
\draw[arr] (c0.north) to[bend left=40] (c5.north);
\node[font=\footnotesize, above] at (3.2, 2.0) {по кругу};

% Метки осей (для отладки)
\foreach \x in {-1,...,8} {
    \node[font=\tiny, gray] at (\x, -1.3) {\x};
}
\foreach \y in {-1,0,1,2,3} {
    \node[font=\tiny, gray] at (-1.3, \y) {\y};
}

\end{tikzpicture}
}% end adjustbox
\caption{DMA circular. Буфер читается по кругу.
Когда DMA доходит до конца~--- начинает сначала.}
\label{fig:dac-dma-circular}
\end{figure}

\begin{keybox}
\textbf{Почему \emph{circular}?} Потому что \emph{сигнал}~---
\emph{периодический}. \emph{Синус} \emph{повторяется} \emph{каждый}
\emph{период}. \emph{Нам} \emph{не} \emph{нужно} \emph{менять}
\emph{таблицу}~--- \emph{достаточно} \emph{читать} её
\emph{по} \emph{кругу}.
\end{keybox}

\subsection{Как это освобождает CPU}

\textbf{Без} DMA:
\begin{enumerate}
    \item \emph{Таймер} \emph{срабатывает} \emph{каждый}
          \emph{такт}.
    \item \emph{Прерывание} \emph{вызывает} \texttt{dds\_tick()}.
    \item CPU \emph{вычисляет} \emph{индекс}, \emph{пишет}
          \emph{в} DAC.
    \item \emph{Возвращается} \emph{к} \emph{работе}.
    \item \emph{Повтор}.
\end{enumerate}

\textbf{С} DMA:
\begin{enumerate}
    \item \emph{CPU} \emph{настраивает} DMA \emph{один} \emph{раз}.
    \item \emph{Таймер} \emph{срабатывает} \emph{каждый}
          \emph{такт}.
    \item \emph{DMA} \emph{читает} \emph{следующее}
          \emph{значение} \emph{из} \emph{таблицы}.
    \item \emph{DMA} \emph{пишет} \emph{в} DAC.
    \item CPU \emph{свободен}.
\end{enumerate}

\begin{keybox}
\textbf{CPU \emph{не} \emph{тратит} \emph{время} \emph{на}
\emph{каждый} \emph{семпл}.} \emph{DMA} \emph{делает}
\emph{это} \emph{сам}. CPU \emph{занимается} \emph{USB},
\emph{командами}, \emph{обработкой}.
\end{keybox}

\section{Калибровка}
\label{sec:dac-calibration}

\textbf{Проблема}: \emph{DAC} \emph{не} \emph{идеален}.
\emph{Есть} \emph{ошибки}:
\begin{itemize}
    \item \textbf{Смещение} (offset): \emph{ноль} \emph{не}
          \emph{точно} \emph{ноль}.
    \item \textbf{Усиление} (gain): \emph{максимум} \emph{не}
          \emph{точно} \emph{максимум}.
    \item \textbf{Нелинейность}: \emph{синус} \emph{не}
          \emph{совсем} \emph{синус}.
\end{itemize}

\textbf{Решение}: \emph{калибровка}. \emph{Измеряем}
\emph{реальный} \emph{выход} \emph{и} \emph{корректируем}
\emph{таблицу}.

\begin{figure}[htbp]
\centering
\adjustbox{max width=\textwidth}{%
\begin{tikzpicture}[
    x=1cm, y=1cm,
    axis/.style={-{Latex[length=3mm]}, thick},
    ideal/.style={very thick, color=clrLogicData},
    real/.style={very thick, color=clrWarn}
]

% ─── Координатная сетка ────────────────────────────
\draw[step=1cm, gray!20, very thin] (-1, -1) grid (9, 6);

% Оси
\draw[axis] (-0.5, 0) -- (8.5, 0) node[right] {\(t\)};
\draw[axis] (0, -0.5) -- (0, 5.5) node[above] {\(U\)};

% Идеальный сигнал
\draw[ideal, domain=0:8, samples=100]
    plot (\x, {2.5 + 1.5 * sin(deg(2 * \x))});

% Реальный сигнал (со смещением и искажением)
\draw[real, domain=0:8, samples=100]
    plot (\x, {2.8 + 1.3 * sin(deg(2 * \x + 0.2))});

% Легенда
\node[font=\footnotesize, fill=white] at (6.5, 5) {идеал};
\node[font=\footnotesize, fill=white] at (6.5, 4.5) {реал};

% Метки осей (для отладки)
\foreach \x in {-1,...,9} {
    \node[font=\tiny, gray] at (\x, -1.3) {\x};
}
\foreach \y in {-1,...,6} {
    \node[font=\tiny, gray] at (-1.3, \y) {\y};
}

\end{tikzpicture}
}% end adjustbox
\caption{Идеальный (зелёный) и реальный (красный) сигналы.
Реальный~--- со \emph{смещением} и \emph{искажением}.}
\label{fig:dac-calibration}
\end{figure}

\subsection{Метод калибровки}

\textbf{Три} \emph{шага}:
\begin{enumerate}
    \item \textbf{Измерить} \emph{выход} \emph{при}
          \emph{известных} \emph{значениях} DAC (\num{0},
          \num{2048}, \num{4095}).
    \item \textbf{Вычислить} \emph{коэффициенты} \emph{смещения}
          \emph{и} \emph{усиления}.
    \item \textbf{Применить} \emph{коррекцию} \emph{к}
          \emph{каждому} \emph{значению} \emph{из} \emph{таблицы}.
\end{enumerate}

\begin{listing}[htbp]
\begin{minted}[frame=none, bgcolor=gray!5]{c}
void dac_calibrate(void)
{
    // 1. Измерить реальные значения
    uint16_t raw_min = adc_read(); // при DAC = 0
    uint16_t raw_max = adc_read(); // при DAC = 4095

    // 2. Вычислить коэффициенты
    float scale = (float)DAC_MAX / (raw_max - raw_min);
    float offset = -raw_min * scale;

    // 3. Применить к таблице
    for (int i = 0; i < TABLE_SIZE; i++) {
        for (int w = 0; w < WAVE_COUNT; w++) {
            uint16_t v = wave_table[w][i];
            v = (uint16_t)(v * scale + offset);
            if (v > DAC_MAX) v = DAC_MAX;
            wave_table[w][i] = v;
        }
    }
}
\end{minted}
\caption{\texttt{dac\_calibrate()}~--- калибровка}
\label{lst:dac-calibrate}
\end{listing}

\begin{warnbox}
\textbf{Калибровка~--- \emph{одноразовая} процедура.}
\emph{Запускается} \emph{при} \emph{старте} (или
\emph{по} \emph{команде}). \emph{Результат}
\emph{сохраняется} \emph{в} \emph{таблице}.
\emph{Дальше}~--- \emph{используется}.
\end{warnbox}

\section{Усилитель}
\label{sec:dac-amplifier}

\textbf{DAC} \emph{выдаёт} \SI{0}{\volt}--\SI{3.3}{\volt}.
\emph{Для} \emph{большинства} \emph{задач} \emph{этого}
\emph{достаточно}. \emph{Но} \emph{иногда} \emph{нужно}
\emph{больше}.

\begin{ideabox}
\textbf{Аналогия: микрофон и усилитель.} \emph{Микрофон}
\emph{выдаёт} \emph{слабый} \emph{сигнал} (\SI{10}{\milli\volt}).
\emph{Усилитель} \emph{увеличивает} \emph{его} \emph{до}
\SI{1}{\volt} \emph{для} \emph{динамиков}. \emph{То} \emph{же}
\emph{с} DAC.
\end{ideabox}

\textbf{Простой} \emph{усилитель} \emph{на} \emph{транзисторе}:
\begin{itemize}
    \item \emph{Коэффициент} \emph{усиления}: $K = R_C / R_E$.
    \item \emph{Для} $K = 2$: $R_C = 2 R_E$.
\end{itemize}

\section{Уроки главы}
\label{sec:dac-dma-lessons}

\begin{enumerate}
    \item \textbf{DMA circular~--- это автопилот.} CPU
          \emph{свободен}.
    \item \textbf{Буфер читается по кругу.} Сигнал
          \emph{периодический}.
    \item \textbf{Калибровка~--- обязательна.} DAC
          \emph{не идеален}.
    \item \textbf{Смещение и усиление~--- две ошибки.}
          \emph{Корректируются} \emph{коэффициентами}.
    \item \textbf{Усилитель~--- если нужно больше.}
          \emph{Простой} \emph{транзистор}.
\end{enumerate}

\section{Чек-лист грабель}
\label{sec:dac-dma-pitfalls}

\begin{table}[htbp]
\centering
\caption{Грабли DMA circular и калибровки}
\label{tab:dac-dma-pitfalls}
\small
\begin{tabularx}{\textwidth}{@{}c L L@{}}
\toprule
\textbf{№} & \textbf{Грабля} & \textbf{Решение} \\
\midrule
1 & DMA не запускается & \texttt{HAL\_DAC\_Start\_DMA} \\
2 & \emph{Нет} сигнала & Проверить \texttt{DMA\_CIRCULAR} \\
3 & Нет калибровки & \texttt{dac\_calibrate()} \\
4 & Переполнение & Ограничить \texttt{DAC\_MAX} \\
5 & Усилитель <<гудит>> & \emph{Добавить} \emph{фильтр} \\
6 & Неточная частота & \emph{Проверить} \emph{таймер} \\
7 & Нет \texttt{volatile} & \emph{Добавить} \\
8 & \emph{Забыли} \texttt{\_\_HAL\_RCC} & Добавить \\
\bottomrule
\end{tabularx}
\end{table}

\section{Упражнения}
\label{sec:dac-dma-exercises}

\begin{exercisebox}
\textbf{Упражнение 1.} Нарисуйте \emph{диаграмму} DMA
circular с \num{4} ячейками. Покажите \emph{позицию}
DMA \emph{через} \num{3} \emph{такта}.
\end{exercisebox}

\begin{exercisebox}
\textbf{Упражнение 2.} Что \emph{произойдёт}, если
\emph{убрать} \texttt{DMA\_CIRCULAR}? \emph{Приведите}
\emph{пример}.
\end{exercisebox}

\needspace{5\baselineskip}
\begin{exercisebox}
\textbf{Упражнение 3.} Реализуйте \emph{калибровку}
\emph{через} \emph{ADC} и \emph{таблицу} \emph{поправок}.
\emph{Как} \emph{сохранить} \emph{результат}?
\end{exercisebox}

\begin{exercisebox}
\textbf{Упражнение 4.} Предложите \emph{схему} \emph{усилителя}
\emph{с} \emph{коэффициентом} \num{3}. \emph{Какие}
\emph{резисторы} \emph{выбрать}?
\end{exercisebox}

\section{Проверьте сами}
\label{sec:dac-dma-check}

\begin{enumerate}
    \item Подключите Nucleo к компьютеру.
    \item Соберите прошивку (см. главу 2).
    \item Подключите осциллограф (HANTEK) к выходу \texttt{PA4}.
    \item Запустите OpenOCD + GDB.
    \item Посмотрите \emph{форму} сигнала на осциллографе.
    \item Запустите \emph{калибровку}. Сравните.
    \item Сделайте \emph{скриншот} — \emph{пришлите} в книгу.
\end{enumerate}

\begin{ideabox}
\textbf{Если получилось~--- вы поняли DMA circular.}
Это \emph{главное}. Дальше~--- \emph{USB}, \emph{ПК},
\emph{GTK4}. \emph{Всё} \emph{впереди}.
\end{ideabox}
